\documentclass{article}\begin{document}The single span considered here has apex the categorical disjoint union of the given apices, and its two functor legs are independently descended from the source and destination legs. At a fixed base object, pack reconstructs its strict source fibre from the disjoint union of the strict source fibres: splitting a vertical arrow uses the actual identity-lift equation, not a choice of preimage up to isomorphism. The equality over identifies the actual destination endpoint functors. WeightedPack then compares totals over those endpoints. Independently, Distribute.gather commutes the contravariant payload Grothendieck total with tagged disjoint union; splitHom recovers the actual tagged grade and the dependent payload arrow. Their composite compare is an equivalence from the tagged sum of weighted fibres to evaluation of the single coproduct span. It commutes as an equality of functors with arbitrary ordinary payload natural transformations, on all arrows, in this forward direction. The proof does not assume that evaluation preserves coproducts. Positive applies this comparison to all positive powers in the previously admitted base and grade universes. All.evaluation includes zero and the positive powers in one actual coproduct span, with the explicit common-size restriction that grade and payload object/hom universes match the base object/hom universes. The zero apex is the original base with both legs identity. The unrestricted family theorem uses a common pair of universe levels among apices; an unrestricted concrete all-length specialization still needs a universe-lifted identity apex or a heterogeneous coproduct with verified compatibility. The fixed-base all-length equivalence composes the inverse equivalence of compare with the summand evaluation equivalences. Payload naturality is proved for compare, not for that composite inverse as a bundled family. Base change still needs the coproduct cleavage with tagged lifts, comparison of its pull to the sum of chosen pulls (including composition and unit), and mixed compatibility. The extension theorem as an equivalence of endofunctors, flattening, free-monad multiplication and strict algebra remain open. See pinned manuscript 458ff36a, lines 929--953; consult Sigma.desc and the definitions of FibreTotal.endpoint and GrothendieckStep.Total to distinguish evaluation of the sum from the sum of evaluations.\end{document}
